\documentclass[aspectratio=169,11pt]{beamer}
\usepackage{teaching_slides}

% Self-study handout: moved out of L1 lecture (Fall 2026).
% Students read this before Lecture 2; ~15 minute summary given in class.

\begin{document}

\title{Econometrics I}
\subtitle{Linear Algebra Review (Self-Study)}
\date{Fall 2026}
\maketitle

\begin{frame}{How to use this handout}
\begin{itemize}
	\item Read this before Lecture 2. It is the linear algebra we rely on all semester --- nothing more.
	\item You will \emph{read} matrix notation constantly in this course; you will almost never manipulate it by hand.
	\item If any of this is new to you, work through the linked visualizations and come to office hours early in the semester.
\end{itemize}
\end{frame}

\begin{frame}{Basic Definitions}
\begin{itemize}
	\item A {\bf vector} in $\mathbb{R}^n$ is a column of numbers $(x_1,x_2,...,x_n)$.
	\item A {\bf matrix} in $\mathbb{R}^{n\times m}$ is $m$ columns of length $n$ vectors (so $n$ is the number of rows and $m$ is the number of columns). We denote an element of a matrix by $m_{ij}$ for row $i$ and column $j$:
	\begin{align*}
		M = \left(\begin{array}{cc}
			m_{11} & m_{12}\\
			m_{21} & m_{22}
		\end{array}\right)
	\end{align*}
	\item For an entity on which there are many pieces of data, we store the data in a vector $x_i$
	\begin{itemize}
		\item[]\emph{Example:} For the USA could have $x_{USA} = \left(GDP_{USA}, Population_{USA},...\right)$
	\end{itemize}
	\item For many entities we can store all the data in a data matrix, $X$.
	\begin{itemize}
		\item[]\emph{Example:} For two countries:
			\begin{align*}
		X = \left(\begin{array}{ll}
			GDP_{USA} & Population_{USA}\\
			GDP_{Canada} & Population_{Canada}
		\end{array}\right)
			\end{align*}
	\end{itemize}
\end{itemize}

\end{frame}


\begin{frame}{Matrix Multiplication}
	\begin{itemize}
		\item For two vectors of equal length define the {\bf dot product} as $v\cdot w$ or $\langle v, w \rangle$:
			\[
				v\cdot w = \sum_{i=1}^n v_i\times w_i
			\]
		\item For two matrices, $A$ and $B$ of sizes $n\times m$ and $m\times k$ define the {\bf matrix product} $C=AB$ as the $n\times k$ matrix with entries $c_{ij}=\sum_{l=1}^ma_{il}b_{lj}$
		\begin{itemize}
			\item Easy way to remember: $(i,j)^{th}$ element of product is dot product of $i^{th}$ row and $j^{th}$ column of $A$ and $B$ respectively.
			\item Not all matrices can be multiplied: left matrix must have column length equal to right matrix's row length
			\item Multiplication is NOT commutative: $AB\ne BA$ even if they both exist
		\end{itemize}
	\end{itemize}
	\url{https://eli.thegreenplace.net/2015/visualizing-matrix-multiplication-as-a-linear-combination/}
\end{frame}


\begin{frame}{Transposes}
\begin{itemize}
	\item Define the {\bf transpose} of $A$ as the matrix $A'$ with elements $a'_{ij} = a_{ji}$ (reverse columns and rows)
	\item A matrix is {\bf symmetric} if $A'=A$
	\item {\bf Important Properties:}
		\begin{itemize}
			\item The matrix $B = A'A$ is \emph{always} a square matrix
			\item The matrix $B=A'A$ is always symmetric
			\item $(A')' = A$
			\item {\bf Multiplication Rule:} $(AB)' = B'A'$
			\item {\bf Addition Rule:} $(A+B)' = A'+B'$
		\end{itemize}
	\item Note that dot products can also be written as an {\bf inner product}: $v\cdot w = v' w$.
	\item Note that the {\bf outer product} of two vectors is a matrix rather than a scalar: $v w'$.

\end{itemize}

\end{frame}

\begin{frame}{Inverses}
\begin{itemize}
	\item The {\bf Identity Matrix}, $I$, is a matrix with 1s on the diagonal and 0s elsewhere. Clearly $AI = A$.
	\item Define the {\bf left inverse} of $A$ to be the matrix $A^{-1}$ such that $A^{-1}A=I$
	\begin{itemize}
		\item Can analogously define right inverse
		\item Right and left inverse will NOT be the same if $A$ is not a square matrix
		\item Right and left inverse WILL be equal if $A$ is square (then we just say inverse)
	\end{itemize}
	\item {\bf Important Properties:}
		\begin{itemize}
			\item {\bf Multiplication Rule:} $(AB)^{-1} = B^{-1}A^{-1}$
			\item {\bf Tranpose Rule:} $(A')^{-1} = (A^{-1})'$
			\item {\bf Dot Product:} $v\cdot w = v'w$
		\end{itemize}
\end{itemize}

\end{frame}

\begin{frame}{Matrix Calculus}
	\begin{itemize}
		\item For a function $f:\mathbb{R}^n\to\mathbb{R}^m$ recall the definition of the derivative or Jacobian of $f$:
			\[
				Df = \left(\begin{array}{ccc}
					\frac{\partial f_1}{\partial x_1} & \dots & \frac{\partial f_m}{x_1}\\
					\vdots & & \\
					\frac{\partial f_1}{\partial x_n} & \dots & \frac{\partial f_m}{x_n}
					 \end{array}\right)
			\]
		\item We WON'T be doing anything too complicated! But we can define two important functions given a vector $x$ and a matrix $A$:
		\begin{itemize}
			\item For $Ax$, $D(Ax) = A$ (as a line in 1-D calc)
			\item For $x'Ax$, $D(x'Ax) = x'(A+A')$ (as a quadratic in 1-D calc)
		\end{itemize}
	\end{itemize}
	The matrix cookbook \url{https://www.math.uwaterloo.ca/~hwolkowi/matrixcookbook.pdf} has a lot (more) helpful properties.
\end{frame}

\begin{frame}{Check yourself}
Before Lecture 2, make sure you can answer:
\begin{enumerate}
	\item If $X$ is $n \times k$, what are the dimensions of $X'X$? Of $X'y$ (with $y$ being $n\times 1$)? Of $(X'X)^{-1}X'y$?
	\item Why is $X'X$ always symmetric?
	\item When does $(X'X)^{-1}$ fail to exist? Give a two-column example.
	\item Expand $(y - X\beta)'(y - X\beta)$ using the transpose rules.
\end{enumerate}
\end{frame}

\end{document}
